%\begin{table*}[h]
%\caption{Architectural and security comparison between Bitcoin and Ethereum in the context of document verification systems.}
%\label{tab:btc_eth_comparison}
%\centering
%\small
%\renewcommand{\arraystretch}{1.18}
%\setlength{\tabcolsep}{4pt}
%
%\begin{tabular}{
%>{\raggedright\arraybackslash\bfseries}p{3.2cm}
%>{\raggedright\arraybackslash}p{6.1cm}
%>{\raggedright\arraybackslash}p{6.1cm}
%}
%\toprule
%Dimension & Bitcoin (PoW, UTXO) & Ethereum (PoS, Account/VM) \\
%\midrule
%
%Security foundation &
%Thermodynamic work: security anchored in continuous energy expenditure; measured by hashrate $H$ &
%Bonded-capital security: security anchored in locked stake; measured by total stake $S$ \\
%
%Influence parameter &
%$\alpha = \dfrac{H_A}{H_{total}}$ &
%$\alpha = \dfrac{S_A}{S_{total}}$ \\
%
%State model &
%UTXO set; transactions consume and produce independent outputs; stateless validation &
%Global account state; transactions mutate shared state; persistent contract storage \\
%
%Execution model &
%Restricted, non-Turing-complete Script; bounded evaluation; no general loops &
%Turing-complete EVM; arbitrary control flow and inter-contract calls \\
%
%Verification execution model &
%On-chain commitment; verification performed off-chain via cryptographic proofs &
%Verification logic typically executed on-chain via smart contracts \\
%
%On-chain cost model &
%Transaction fee proportional to virtual bytes; no execution gas model &
%Gas-based execution cost; fees proportional to computation and storage \\
%
%Smart contract dependency &
%No native smart contract state; minimal script functionality &
%Extensive reliance on smart contracts for certification and attestation logic \\
%
%Revocation handling &
%Requires commitment update via new transactions &
%Implemented via contract state mutation \\
%
%Consensus finality &
%Probabilistic finality (Nakamoto-style) &
%Economic finality after checkpointing \\
%
%Base-layer attack surface &
%Minimal execution engine; reduced systemic complexity &
%Larger execution surface including VM runtime and contract standards \\
%
%Protocol governance model &
%Conservative soft-fork upgrades (e.g., SegWit, Taproot) &
%Multiple major hard-fork transitions (e.g., DAO fork, Merge) \\
%
%\bottomrule
%\end{tabular}
%\end{table*}



\begin{table*}[t]
\caption{Condensed architectural and security comparison between Bitcoin and Ethereum for document verification systems.}
\label{tab:btc_eth_comparison_short}
\centering
\small
\renewcommand{\arraystretch}{1.15}
\setlength{\tabcolsep}{4pt}

\begin{tabular}{
>{\raggedright\arraybackslash\bfseries}p{3.0cm}
>{\raggedright\arraybackslash}p{6.2cm}
>{\raggedright\arraybackslash}p{6.2cm}
}
\toprule
Dimension & Bitcoin (PoW, UTXO) & Ethereum (PoS, Account/VM) \\
\midrule

Security basis &
Energy-based security (hashrate $H$) &
Stake-based security (total stake $S$) \\

Adversarial threshold &
$\alpha = H_A/H_{total}$ &
$\alpha = S_A/S_{total}$ \\

State model &
UTXO; stateless validation &
Global account state; persistent storage \\

Execution model &
Limited, non-Turing Script &
Turing-complete EVM \\

Verification approach &
On-chain commitment, off-chain proof verification &
On-chain contract-based verification \\

Cost model &
Fee per transaction size (vbytes) &
Gas-based fee (computation + storage) \\

Revocation mechanism &
New commitment transaction required &
Contract state update \\

Finality model &
Probabilistic finality &
Economic finality \\

Attack surface &
Minimal scripting layer &
VM and smart-contract dependent \\

Governance style &
Conservative soft-forks &
Frequent hard-fork upgrades \\

\bottomrule
\end{tabular}
\end{table*}
